\cotefile{136}
\scanpage{2}
Complexe DR à p.d. et ombres

% archive sheet number: 1
\scanpage{3}
\mspage{1}
Complexe DR à p.d. et ombres des modèles.

% slanted
\margin{Plan conf. IHES en 1976}

1 Historique

a) Notion des formes diff. (\uline{Poincaré} ?) et formule de \uline{Stokes}
\[
\int_c d\omega=\int_{\partial c}\omega
\qquad\text{(d'où}\quad H^{\ast}_{\mathrm{DR}}(X)\longrightarrow H^{\ast}(X,\mathbb C)
\]
\margin{géom. diff\\ -- anal.\\ -- algébr.\\ \deleted{\unsure{top.}} arithm.\\ top. alg $\left\{\begin{array}{l}\mathrm{PL}\\1/2\,\unsure{\mathrm{simpl.}}\end{array}\right.$}

b) Th. de \uline{De Rham}. (Conjecturé par \uline{É. Cartan}). \unsure{Mais du point de présentation|Mais de présentation|Mais on perd la torsion}

Maintenant bien compris : th. des faisceaux, lemme Poincaré

c) Théorie de \uline{Hodge des} intégrales harmoniques (structure supplémentaire bigraduée sur $H^{\ast}_{\mathrm{DR}}(X)$ si $X$ kählérienne compacte)

d) Théorème de \uline{H. Cartan}-\uline{Serre} sur variétés de Stein (car lemme de Poincaré valable dans l'analytique)

e) Cas des variétés algébriques ou schémas sur \inserted{corps de} base (ou schéma de base général) : \uline{Dwork}, \uline{Washnitzer}-\uline{Monsky}, plus tard le yoga ``cristallin'' développé par \uline{Berthelot}, \uline{Illusie}, \deleted{\unsure{et repris par d'autres}}\inserted{\uline{Messing}, Mazur}
% slanted
\margin{cf aussi Hartshorne Herrera Ogus, Bloch (?)}

Ici on trouve des théories coh. qui ne sont plus à ``anneau de coeff''\unsure{\}} de car. nulle i.e contenant $\mathbb Q$ -- i.e. en principe les phénomènes de torsion ne sont pas perdus. \inserted{Du pt de vue Weil, \unsure{s'était|c'était} un défaut impardonnable !}

f) Th. de \uline{Grothendieck} pour variétés algébriques sur $\mathbb C$ (\deleted{\unsure{affiné}}\inserted{\uline{généralisé}} par \uline{Deligne}, \uline{Hartshorne} pour des coeff plus généraux) Ceci donne confiance (du moins en car. 0) en la signification topologique de la coh. de De Rham ``algébrique'' des schémas algébriques.

g) Complexe de De-Rham-Sullivan pour espaces topologiques généraux : \uline{formes différentielles singulières} $C^\infty$ \deleted{\illegible{2} $\unsure{C_{\mathrm{alg}}}$ \unsure{algébrique, resp.}} ($\mathbb C$-algébrique, resp. $\mathbb R$-algébrique, resp $\mathbb Q$-algébrique). Donne la coh. à coeff dans $\mathbb C$ (resp. $\mathbb R$, resp. $\mathbb Q$) (facile) Sullivan montre mieux.

\diagram{\begin{tikzcd}[ampersand replacement=\&,column sep=small,row sep=small]
C^\bullet_{\mathrm{DRS}\,\mathbb R\text{-}\mathrm{alg}} \arrow[r,phantom,"\subset"] \& C^\bullet_{\mathrm{DRS}\,C^\infty} \arrow[r,phantom,"\subset"] \arrow[d,phantom,"\supset"{rotate=-90}] \& C^\bullet_{\mathrm{DRS}} \\
\deleted{\unsure{C}} \& C^\bullet_{\mathrm{DR}} \&
\end{tikzcd}}
% archive sheet number: 2
\scanpage{4}
\mspage{2}
que le $\mathbb Q$-type d'homotopie de $X$ est récupéré si $X$ est simplement connexe -- de façon plus précise, il y a \inserted{(sauf erreur)} une équivalence de cat. \inserted{donnée par $C^{\ast}_{\mathrm{DRS}}$} entre la cat. hom. \inserted{faible} des espaces connexes et simplement connexes (du pt de vue singulier), et une certaine catégorie dérivée formée avec les $\mathbb Q$-algèbres graduées associatives et \inserted{anti}comm. \deleted{telles} à degrés $\geq 0$ telles que
\[
H^0(A)\simeq\mathbb Q\ \deleted{\unsure{H^1(A)\simeq 0}}\quad H^1(A)\simeq 0 .
\]
Il y a une théorie des modèles minimaux \inserted{pour de telles algèbres}, une façon très simple de récupérer les $\pi_i\otimes_{\mathbb Z}\mathbb Q$ en termes d'un tel modèle.--- (On renvoie au papier de Sullivan)

Mais à nouveau, on perd la torsion !

\noindent\rule{\linewidth}{0.3pt}

Dém. du th. d'isomorphie de Sullivan
\[
H^\bullet(C^\bullet_{\mathrm{DRS}\deleted{\unsure{\mathbb Q}}}(X,\mathbb Q))\simeq H^\bullet(X,\mathbb Q)
\]

$\forall n\geq 0$, $[n]$ simplexe type de dim $n$
\[
\left(\sum_0^n X_i=1,\ \deleted{\unsure{\varepsilon}}X_i\geq 0\right)
\]
$\Delta^{[n]}\deleted{\unsure{\mathbb R}}$ sa réalisation géométrique

$E^{[n]}$ l'espace affine qui le sous tend (avec une $\mathbb Q$-structure)
\margin{$\displaystyle\left|\begin{array}{l}\sum_0^n X_i=1\\ \text{dans }\mathbb R^{n+1}\end{array}\right.$}

\[
\mathrm{DRS}^{\bullet}_{\deleted{\unsure{i}}[n]}=C^\bullet_{\mathbb Q\text{-}\mathrm{DR}}(E^{[n]})\quad\text{son complexe de De Rham $\mathbb Q$-algébrique}
\]
C'est contravariant en $[n]$, d'où
\[
\mathrm{DRS}^{\bullet\deleted{\unsure{\ast}}}_{\ast}=(\mathrm{DRS}_{[n]})_{n\geq 0}
\]
Algèbre différentielle graduée \deleted{(semi)}\inserted{semi-}simpliciale (et $=$ simplicial) -- à degrés $\geq 0$, anticommutative.

Pour tt espace topologique $X$, $S_{\ast}(X)$ son complexe singulier $1/2$ simplicial. On a
\[
C^\bullet_{\mathrm{DRS}}(X,\mathbb Q)\simeq\mathrm{Hom}(S_{\unsure{\ast}}(X),\mathrm{DRS}^{\bullet}_{\ast})
\]

% archive sheet number: 3
\scanpage{5}
\mspage{3}
i.e. \quad a) $C^\bullet_{\mathrm{DRS}}(X,\mathbb Q)$ dépend de $X$ via $S_{\unsure{\ast}}(X)$

b) Sur $(\mathrm{Ss})\deleted{\unsure{\ast}}=(\mathrm{Ens})_{\ast}$, le foncteur $C^\bullet_{\mathrm{DRS}}(-,\mathbb Q)$ est représentable par $\mathrm{DRS}^{\bullet}_{\ast}$.

Or \quad a) $\mathrm{DRS}^{\bullet}_{\ast}$ est une résolution de $\mathbb Q_{\ast}$ \inserted{(dans la cat. des groupes $1/2$ simpl.)} (lemme de Poincaré algébrique sur les espaces \inserted{$\mathbb Q$-}affines $E^{[n]}$)

b) Les composantes $\mathrm{DRS}^{i}_{\ast}$ $(i\geq 0)$ de $\mathrm{DRS}^{\bullet}_{\ast}$ sont des objets \inserted{\unsure{ab}} acycliques du topos $\mathrm{Ss}$ (ce qui revient à dire que leurs $\pi_j(\mathrm{DRS}^{i}_{\ast})$ sont nuls, ce qui se vérifie facilement. \inserted{Cela serait itou avec des formes $C^\infty$ à coeff. ds $\mathbb R$ ou $\mathbb C$, ou analytiques réelles, \deleted{ou} analytiques complexes, \ldots)}

\deleted{\unsure{On peut partir du}} On aimerait avoir une $\mathbb Z$-algèbre diff. $C_{\mathrm{DR}}(X,\mathbb Z)$ \uline{alternée} (dépendant fonct. de $X$, ou $S_{\ast}(X)$) dont la coh. soit la coh. entière de $\mathbb Z$. Mais il paraît qu'il y a des ths contre ! On va pourtant y arriver \deleted{\unsure{quand}}\unsure{en un sens|en q. sens} !

Si on prend $C_{\mathbb Z\text{-}\mathrm{DR}}^{\bullet}(E^{[n]})$ (car $E^{[n]}$ a une $\mathbb Z$-structure affine) c'est a) qui devient déjà faux : pour intégrer $\int x^n dx$, il faut un dénominateur \inserted{avec} $\frac{x^{n+1}}{n+1}$ ! Mais en géom. alg. on est déjà familiarisé avec une façon de sauter à pieds-joints \deleted{sur}\inserted{par dessus} le caniveau, en introduisant des p.d. et des polynômes \inserted{(ou séries formelles)} à puiss. div. Si $E^{[n]}$ avait une origine sur $\mathbb Z$ (i.e. provenait can. d'un $\mathbb Z$-module libre de déf.) on aurait un complexe de De Rham à p.d. Mais pas pour un espace affine !. Notons
\[
\begin{aligned}
\mathrm{DRS}^{\bullet}_{[n]}&\simeq\deleted{\mathbb Q[X_i,dX_i]/}\ C^\bullet_{\mathrm{DR}/\mathbb Q}\big(\mathbb Q[(X_i)_{0\leq i\leq n}]/\textstyle\sum X_i-1\big)\\
&\simeq\mathbb Q[X_i,dX_i]_{0\leq i\leq n}/\big(\textstyle\sum X_i-1,\ \textstyle\sum dX_i\big)\\
&\simeq\mathbb Q[X_0,\ldots,X_{n-1}][dX_0,\ldots,dX_{n-1}]
\end{aligned}
\]

% archive sheet number: 4
\scanpage{6}
\mspage{4}
Donc on aurait envie de prendre
\[
\deleted{\mathbb Q}\ \mathbb Z\big(\{\unsure{(}X_i\unsure{)}\}[dX_i]\big)_{0\leq i\leq n}/\big(\textstyle\sum X_i-1,\sum dX_i\big)_{\deleted{\mathrm{pd}}}
\]
où $\{\ \}$ désigne les pol. à p.d., mais ce n'est plus $\simeq\mathbb Z\{X_0,\ldots,X_{n-1}\}[dX_0,\ldots,dX_{n-1}]$ \deleted{(qui est} \deleted{\unsure{le est|la cat}}\inserted{bien} une résolution de $\mathbb Z$), car si on veut tirer de la relation $\sum_0^n X_i-1=0$,
\[
X_n=1-\sum_0^{n-1}X_i\quad(\inserted{\text{et}}\ dX_n=-\sum_0^{n-1}dX_i\quad\text{de }\sum_{\unsure{1|0}}^n dX_i=0)\ .\text{--}
\]
on a le \unsure{``bec''|``hic''} que $1-\sum_0^{n-1}X_i$ n'appartient pas à l'idéal à p.d. donné dans $\mathbb Z\{X_0,\ldots,X_{n-1}\}$ (donc on ne voit pas comment envoyer $\mathbb Z\{X_0,\ldots,X_{n-1}\}$ dans $\mathbb Z\{X_0,\ldots,X_{n-1}\}$ avec \unsure{comme|même} \deleted{noyau} $\sum_{\unsure{1|0}}^n X_i-1$ dans le noyau\ldots).

On s'en tire en prenant un anneau de coeff. différent de $\mathbb Z$, soit $S$, avec un \deleted{indéterminé}\inserted{``paramètre''} $t\in S$ fixé dont on \uline{sache prendre les p.d.} \inserted{(i.e. $t\in J$, $J$ idéal à p.d. de $S$)}, et en remplaçant les équations $\sum X_i=1$ de $E^{[n]}$ par l'équation
\[
\sum X_i=t\qquad\text{dans } S^{n+1}
\]
et définissant
\[
\big(C^\bullet_{\mathrm{DRpd}\,[n]}(S,J,t)\big)\unsure{\simeq|=}
\big(S\{X_i\}_{0\leq i\leq n}[dX_i]_{0\leq i\leq n}\big)/
\big(\textstyle\sum X_i-t,\sum dX_i\big)
\]
(On divise par l'idéal \uline{à p.d.} engendré \unsure{pd} par $\sum X_i-t,\ \sum dX_i$, ce qui a un sens
% archive sheet number: 5
\scanpage{7}
\mspage{5}
car dans $S\{X_i\}[dX_i]$ l'idéal formé des formes \inserted{à p.d.} d'augmentation $\in J$ est à p.d.
\[
\big(S\{X_i\}_{0\leq i\leq n}/(\textstyle\sum X_i-t)_{\mathrm{pd}}\big)
\otimes_S\bigwedge^{\ast}\big(S^{[n]}/\mathrm{diag}\,S^{[n]}\big)
\]
C'est une $S$-algèbre diff. \inserted{graduée} anti comm., à degrés $\geq 0$, \inserted{augmentée vers $S/J$, et} à p.d. sur l'idéal \deleted{engendré par $\bigwedge^+,S\{X_i\}^+$ et $J\{X_i\}^0$ \unsure{isomorphe au}} noyau de l'augmentation
\[
C^{\bullet}_{\mathrm{DRpd}\,[n]}(S,J,t)\longrightarrow S/J\ ;
\]
% slanted
\margin{la diff est compatible avec les p.d. i.e. $d(x^{(n)})=x^{(n-1)}dx$}
et comme telle isomorphe à $S\{X_0,\ldots,X_{n-1}\}[dX_0,\ldots,dX_{n-1}\unsure{\}|\mathclose{]}}$ qui \uline{est une résolution de $S$}. \deleted{Donc a) est acquis. Mais faute de b), on \inserted{ne} peut dire en général ce \deleted{qu'il en soit} qu'est $H^{\ast}C^\ast$.}

Pour $[n]$ variable, on trouve
\[
C^{\bullet}_{\mathrm{DRpd}\,\ast}(S,J,t)=
\big(C^{\bullet}_{\mathrm{DRpd}\,[n]}(S,J,t)\big)_{n\geq 0}
\]
qui est une rés. $1/2$ simpliciale (et $=$ simpliciale) de $S$, \inserted{avec aug. vers $S/J=k$, et p.d. sur l'idéal d'augm., \unsure{compatible}} Elle dépend fonctoriellement de $(S,J,t)$, et elle permet de former, \inserted{pour $\unsure{\mathscr X|X|\mathcal X}_{\ast}\in(\mathrm{Ss})$}
\[
C^{\bullet}_{\mathrm{DRpd}}(\mathscr X_{\ast};S,J,t)=\mathrm{Hom}\big(\mathscr X_{\ast},C^{\bullet}_{\mathrm{DRpd}\,\ast}(S,J,t)\big)
\]
foncteur contravariant en $\mathscr X_{\ast}$\unsure{!} (et si $X$ esp. top
\[
\begin{aligned}
C^{\bullet}_{\mathrm{DRpd}}(X;\qquad)&=C^{\bullet}_{\mathrm{DRpd}}(S_{\ast}(X);\text{---})\\
&=\mathrm{Hom}\big(S_{\ast}(X),C^{\bullet}_{\mathrm{DRpd}\,\ast}(\quad)\big).
\end{aligned}
\]
% margin note written sideways
\margin{\unsure{!} à valeurs dans les $\unsure{S|\mathbb Z}$-algèbres graduées diff. $S/J$-augmentées à p.d. sur l'idéal d'augmentation, \unsure{compatibles|comp. tibles} avec la diff.}

Mais on ne peut dire en général \inserted{(on a seulement $H^{\bullet}_{\mathrm{DRpd}}(\mathscr X_{\ast};S,J,t)\longrightarrow H^\bullet(\mathscr X_{\ast},S)$)} quelle est sa cohomologie, et en tous cas elle ne sera pas isom (pour $\mathscr X_{\ast}$ variable) à $H^\bullet(\mathscr X_{\ast},S)$.

% archive sheet number: 6
\scanpage{8}
\mspage{6}
\inserted{Mais} Soit $k$ un anneau comm (ass. unit.) et $T$ une \unsure{indéterminé|indéterminée}, on prendra dorénavant
\[
\begin{aligned}
S&=k\{T\},\qquad J=k\{T\}^{+}=\operatorname{Ker}(k\{T\}\unsure{\text{---}|\to} k)\\
t&=T
\end{aligned}
\]
Donc
\[
\left\{\begin{aligned}
C^\bullet_{\mathrm{DRpd}\,[n]}(S,J,t)&\simeq
\underbrace{k\{T,X_0,\ldots,X_n\}/(\textstyle\sum X_i-T)_{\mathrm{pd}}}_{\substack{\wr\\ k\{X_0,\ldots,X_n\}}}
\otimes_k\bigwedge^{\bullet}k^{[n]}/k\\
S/J&\simeq k
\end{aligned}\right.
\]
Soit
\[
\begin{aligned}
\Phi_{k\ast}&=\big[[n]\longmapsto k^{[n]}\big]\quad\underset{\substack{\text{inclusion}\\\text{diagonale}}}{\supset}\ k_\ast=\big[[n]\text{---}k\big]\\
\Psi_{k\ast}&=\Phi_{k\ast}/\unsure{(k_\ast)|k_\ast}=\big([n]\text{---}k^{[n]}/\underbrace{k}_{\mathrm{diag}}\big)
\end{aligned}
\]
On a alors
\[
\begin{aligned}
C^\bullet_{\mathrm{DRpd}\,\ast\deleted{\unsure{n}}}(k\{T\},k\{T\}^{+},k)&\simeq\unsure{\Gamma|\mathscr T}^{\bullet}_{k}\Phi_{k\ast}\otimes_k\bigwedge^\bullet\Psi_{k\ast}\\
\underset{\text{déf}}{\Vert}&\\[-2pt]
C^{\bullet\bullet}_{\mathrm{DRpd}\,\ast}(k)&
\end{aligned}
\]
\[
\boxed{C^{\bullet\bullet}_{\mathrm{DRpd}\,\ast}(k)\simeq\Gamma^\bullet_k\Phi_{k\ast}\otimes_k\bigwedge^\bullet\Psi_{k\ast}}
\]

\uline{Structure}. \inserted{$k\{T\}$-}Algèbre \deleted{\unsure{diff.}} \uline{bigraduée} (deg. \unsure{complém.} \uline{et} degré extérieur, d'\uline{où} degré total) \inserted{unit.} assoc. \deleted{anticomm.}\inserted{alternée} (anticomm + carré des éléments de degré impair nuls), augmentée vers $k$, à p.d. sur l'idéal d'augm, avec diff de bidegré $-1,+1$ (donc degré total 0) \inserted{compatible avec p.d.}
% slanted
\margin{la bigraduation provient de la \uline{graduation} de $S$, et du fait que $t$ est homogène (ici de degré 1)}

% archive sheet number: 7
\scanpage{9}
\mspage{7}
Ces structures sont héritées par les
\[
\boxed{\begin{aligned}
C^{\bullet\bullet}_{\mathrm{DRpd}}(\unsure{\mathscr X|X|\mathcal X}_{\ast},k)&\overset{\mathrm{déf}}{=}C^\bullet_{\mathrm{DRpd}}(\mathscr X_{\ast};k\{T\},k\{T\}^{+},T)\\
&=\mathrm{Hom}\big(\mathscr X_\ast,C^{\bullet\bullet}_{\mathrm{DRpd}\,\ast}(k)\big)
\end{aligned}}
\]
et dépendent de façon contravariante de $\mathscr X_{\ast}$, (covariante de $k$).

$\Phi_{k\ast}$ est \deleted{\unsure{$k$}}\inserted{un $k$-Module $1/2$ simplicial} homotope à 0 donc $\unsure{\Gamma|\mathscr T}^{p}_k(\Phi_{k\ast})\otimes\bigwedge^q\Psi_{k\ast}$ \deleted{\unsure{est}} est homotope à $\Gamma^p_k(0)\otimes\bigwedge^q\Psi_{k\ast}$, donc hom. à 0 si $p\ne 0$.

Donc $C^{p,q}_{\mathrm{DRpd}\,\ast}(k)$ est hom à 0 (donc \deleted{flasque}\inserted{acyclique}) si $p=0$. En degré total donné \inserted{$n$}, on a
\diagram{\begin{tikzcd}[ampersand replacement=\&,column sep=tiny,row sep=small]
0 \arrow[r] \& C^{n,0}_{\ast} \arrow[r] \& C^{n-1,1}_{\ast} \arrow[r] \& C^{n-2,2}_{\ast} \arrow[r,dash] \& \cdots \& C^{\unsure{1|0},n-1}_{\ast} \arrow[r,phantom,"\unsure{\text{---}|\leftarrow|\to}"] \& C^{0,n}_{\ast} \arrow[r] \& 0 \arrow[d] \\
\& k_\ast \arrow[u] \& \& \& \& \& \& 0 \arrow[d] \\
\& \& \& \& \& \& \& 0
\end{tikzcd}}
i.e. on trouve une résolution \inserted{de long. $n$} \deleted{flasque}\inserted{de} $k_{\ast}$ par des \deleted{groupes}\inserted{$k$-Modules} $1/2$ simpliciaux qui sont acycliques sauf le dernier -- donc on peut la considérer comme \unsure{une|un} tronqué de \deleted{l'indice}\inserted{degré} $n$ d'une résolution flasque de $k_{\ast}$ -- la cohomologie de ses sections sur un $\mathscr X_{\ast}$ est donc la coh. de $\mathscr X_{\ast}$ tronquée en degré $n$ :
\[
H_{\mathrm{DRpd}}^{p,q}(\mathscr X_{\ast},k)\simeq
\begin{cases}H^q(\mathscr X_{\ast},k)&\text{si }q\leq p+q\ \text{i.e. }p\geq0\\
0&\text{si }p<0
\end{cases}
\]

% archive sheet number: 8
\scanpage{10}
\mspage{8}
Mais quelle est \inserted{(pour $q$ fixé)} la structure de $k\{T\}$-module de $H^{\bullet\,q}_{\mathrm{DRpd}}(\unsure{\mathscr X|X|\mathcal X}_{\ast},k)$ ? \deleted{\unsure{On calcule ce}}\inserted{On voit que} pour le \uline{degré total} (égal au degré extérieur. \uline{plus} $q$) on trouve $H^0(\mathscr X_\ast,k)\otimes_k k\{T\}$ tronqué en degré $\geq q$, donc
\[
H^{\bullet\,q}_{\mathrm{DRpd}}(\mathscr X_{\ast},k)\simeq\tau_q\big(H^0(\mathscr X_{\ast},k)\otimes_k k\{T\}\big)[q]
\]
\margin{(translation des degrés par $-q$)}

Si on réindexe les \unsure{bidegrés} par le couple (degré total, degré extérieur)
\[
{H'}^{n,q}_{\mathrm{DRpd}}(\mathscr X_{\ast},k)=H^{\overbrace{n-q}^{p},q}_{\mathrm{DRpd}}(\mathscr X_{\ast},k)
\]
(donc la condition de degré $p,q\geq0$ devient $n\geq q\geq 0$, l'op. diff est de bidegré $(0,1)$ donc c'est un hom (\uline{homogène}) de $S$-modules gradués), on trouve
\[
{H'}^{\bullet,q}_{\mathrm{DRpd}}(\mathscr X_\ast,k)\simeq\tau_q\big(H^q(\mathscr X_\ast,k)\otimes_k k\{T\}\big)
\]
Ces isom. sont compatibles avec les structures multiplicatives (et bien entendu fonctoriels en $\mathscr X_{\ast},k$ \ldots). Donc \inserted{\uline{a priori}} on ne récupère \inserted{(via $\mathord{H'}^{\bullet\bullet}_{\mathrm{DRpd}}(\mathscr X_{\ast},k)$)} pas les $k$-modules $H^q(\mathscr X_{\ast},k)$ et leurs \inserted{cup-}accouplements, mais les $k$-modules \inserted{gradués} $\tau_q(H^q(\mathscr X_{\ast},k)\otimes_k k\{T\})$ et leurs accouplements déduits du cup des $H^\bullet(\mathscr X_{\ast},k)$,
% archive sheet number: 9
\scanpage{11}
\mspage{9}
Ce qui donne un espoir que le complexe de DR à p.d. de $\unsure{\mathscr X|X|\mathcal X}_\ast$ \inserted{à coeff ds $k=\mathbb Z$ disons} (même dans le cas $k=\mathbb Q$ \unsure{\}|)} qu'il contient) donne \unsure{une|bien} information homotopique précise sur $\mathscr X_\ast$, c'est l'observation qu'en fait, pour un \deleted{module gradué} $k$-module $M$ quelconque, on récupère $M$ à isom canonique près par la connaissance d'un quelconque des tronqués $\tau_q(M\otimes_k S^\bullet)$ (quelque grand que soit $q$\ldots) \inserted{\unsure{que j'appelle|qu'on appelle} la \uline{$q$ième ombre de $M$}.}, et de $=$ tt accouplement $M\otimes N\unsure{\text{---}|\to}P$ est connu quand on connaît, pour $q$ assez grand, l'accouplement correspondant $\tau_q(M_S)\otimes\tau_q(N_{\unsure{S|J}})\unsure{\text{---}|\to}\tau_q(P_S)$. Donc une \inserted{$k$-}algèbre graduée $H^\bullet$ à degrés $\geq 0$ \deleted{est} est connue à isom canonique près quand on connaît la $S$-algèbre bigraduée dont les composantes \inserted{de degré} ``extérieur'' \inserted{$q$} sont les $\tau_q(H^q\otimes_k S)$ \ldots. \deleted{Ainsi} Cette ``théorie des ombres'' étant \inserted{supposée} acquise, on voit que la \uline{connaissance} de $C^{\bullet\bullet}_{\mathrm{DRpd}}(\mathscr X_\ast,k)$ \inserted{(en tant que $S$-algèbre diff. bigraduée)} implique celle de l'algèbre graduée $H^{\unsure{\ast|\bullet}}(\mathscr X_\ast,k)$ -- et en raffinant un peu, on \inserted{voit qu'on} trouve \unsure{que|$=$} $R^{\bullet}\unsure{\Gamma|\mathscr T}(\mathscr X_\ast,k)$ comme élément de la cat. dérivée $D(k)$, \uline{plus} l'accouple%
% archive sheet number: 10
\scanpage{12}%
\mspage{10}%
ment de cup-produit
\[
\mathbb R\Gamma(\unsure{\mathscr X|X|\mathcal X}_{\ast},k)\overset{L}{\otimes}\mathbb R\Gamma(\mathscr X_\ast,k)\unsure{\text{---}|\to}\mathbb R\Gamma(\mathscr X_\ast,k)
\]

\uline{Remarque et Pbs}

a) Si $k$ est une $\mathbb Q$-algèbre, de sorte que $C^\bullet_{\mathrm{DRS}}(\mathscr X_{\ast},k)$ est défini, on le reconstruit à partir de $C^{\bullet\bullet}_{\mathrm{DRpd}}(\mathscr X_{\ast},k)$ \deleted{\unsure{comme}} par la formule
\[
C^\bullet_{\mathrm{DRS}}(\mathscr X_\ast,k)\simeq C^{\bullet\bullet}_{\mathrm{DRpd}}(\mathscr X_\ast,k)/(T-1)
\]
Plus \unsure{généralemt|généralement}, \uline{quel que soit} $k$, on a
\[
C^\bullet_{\mathrm{DRS}}(\mathscr X_\ast,\underbrace{k\otimes_{\mathbb Z}\mathbb Q}_{k_{\mathbb Q}})\simeq C^{\bullet\bullet}_{\mathrm{DRpd}}(\mathscr X_\ast,k)/(T-1)
\]

[et plus gén. encore, si $(S,J,t)$ comme au début
\margin{\unsure{??|?}}
\[
C^\bullet_{\mathrm{DRpd}}(\mathscr X_{\ast};S,J,t)/(t-1)\simeq C^\bullet_{\mathrm{DRS}}(\mathscr X_\ast,S_{\mathbb Q}/\unsure{k} (t-1))\big]
\]

b) Je suis convaincu que la structure à p.d. sur l'idéal d'augmentation de
\[
\unsure{C|kC}^{\bullet\bullet}_{\mathrm{DRpd}}(\mathscr X_{\ast},k)\unsure{\text{---}|\to}H^{\unsure{0|q}}(\mathscr X_\ast,k)=\mathrm{Hom}(\mathscr X_{\ast},k_\ast)
\]
est importante. \deleted{\uline{Il n'est pas exclu que}} \unsure{l'idéal|l'idée} de se poser la question, pour $\mathscr X_\ast$ \unsure{simplemt|simplement} connexe, si la connaissance de $C^{\bullet\bullet}_{\mathrm{DRpd}}(\mathscr X_\ast,k)$ \uline{avec ttes ces structures} (y compris celle des p.d.) n'implique pas la connaissance du type d'homotopie de $\mathscr X_\ast$. On peut
% slanted
\margin{avec conditions de finitude sur $\mathscr X_\ast$ savoir les $H_i(\mathscr X_\ast)$ de t.f.}
% archive sheet number: 11
\scanpage{13}
\mspage{11}
plus précisément appeler \uline{complexe de De Rham à pd virtuel} \inserted{\uline{sur $k$}} une \deleted{complexe de} $k$-bi-algèbre diff. \deleted{\unsure{k}}augmentée, ass. unit. alternée, à diff de bidegré $-1,+1$, \deleted{(ou $0,+1$ au choix)} à bidegrés $\geq 0$, avec p.d. sur l'idéal d'augm, compatibles avec le diff ($d(x^{(n)})=x^{(n-1)}dx$), et telle que les $H^{\bullet q}(C^{\bullet\bullet})[-q]$ \deleted{sont} des $q$-ombres \unsure{$\simeq$|$\sim$} (auquel cas on récupère à partir de $C^{\bullet\bullet}$ un élément de $D(k)$ avec structure multiplicative ass. unit. comm. $\ldots$), passer à une ``catégorie dérivée'' de ces complexes en inversant les flèches qui sont des \unsure{qu-isom|quasi-isom|q-isom.}, d'où par $C^{\bullet\bullet}_{\mathrm{DRpd}}(\unsure{\mathscr X|X|\mathcal X},k)$ un foncteur
\diagram{\begin{tikzcd}[ampersand replacement=\&]
(\mathrm{Hot})\arrow[r] \& (\mathrm{DRpd})\\[-2pt]
\text{types d'hom} \& \substack{\text{cat. dérivée}\\\text{des complexes}\\\mathrm{DR\ pd}}
\end{tikzcd}}
et on peut se demander si sa restriction aux espaces connexes et simplemt connexes avec des $H_{\deleted{\unsure{\ast|i}}}^{i}(X,\mathbb Z)$ de t.f. (cas $k=\mathbb Z$) \deleted{est} \inserted{induit} une équivalence avec les complexes de DR à pd sur $\mathbb Z$ tels que $H^0(C^{\bullet\bullet})\simeq k$, $H^1(C^{\bullet\bullet})=0$, les $H^i(C^{\bullet\bullet})$ de t.f. sur $\mathbb Z$ \ldots

% archive sheet number: 12
\scanpage{14}
\mspage{12}
c) J'ai un peu réfléchi \deleted{\unsure{aux}} si on peut réconstruire les op. coh. \deleted{\unsure{dans}} (type Steenrod ou Whitney) dans la coh. de $\unsure{\mathscr X|X|\mathcal X}_\ast$ par la connaissance de $C^{\bullet\bullet}_{\mathrm{DRpd}}(\mathscr X_\ast,-)$, et n'ai que des résultats \inserted{partiels} négatifs qui montrent qu'avec des \unsure{autom|automs|autres} multiplicatifs de $C^{\bullet\bullet}_{\mathrm{DRpd}\,\ast}(\deleted{\unsure{\mathbb Z}}k)$ on ne trouve rien d'intéressant.

d) Il faudrait sans doute chercher des modèles minimaux à la Sullivan, pour essayer entre autres d'exprimer les \deleted{\unsure{t}} groupes d'homotopie de $\unsure{X|\mathscr X|\mathcal X}_\ast$ en termes de $C^{\bullet\bullet}_{\mathrm{DRpd}\,\ast}(\mathscr X_\ast,\mathbb Z)$ (cas où $\mathscr X_\ast$ \deleted{\unsure{simplemt|simpl.}} connexe avec cond. de finitudes \ldots) Je n'ai rien fait dans cette direction, Je n'ai $=$ pas développé une formule de Künneth pour $C^{\bullet\bullet}_{\mathrm{DRpd}}$ d'un produit $\mathscr X_\ast\times\mathscr Y_\ast$ -- il y a des difficultés techniques dues au fait qu'on ne peut sans doute supposer $C^{\bullet\bullet}_{\mathrm{DRpd}}(\mathscr X_\ast\ \text{ou}\ \mathscr Y_\ast,k)$ plats sur $k\{T\}$.

% archive sheet number: 13
\scanpage{15}
\mspage{13}
e) Le complexe de chaînes $C_\bullet(\unsure{\mathscr X|X|\mathcal X}_\ast,\deleted{\unsure{\mathbb Z|k}}k)$ permet de réconstruire tous les complexes de cochaînes $C^\bullet(\mathscr X_\ast,k')$, pour $k'$ \deleted{variable} \inserted{$k$-alg.}, par
\[
C^\bullet(\mathscr X_\ast,k)\simeq\mathrm{Hom}_{\mathbb Z}(C_\bullet(\mathscr X_\ast,\deleted{\unsure{\mathbb Z}}k),k).
\]
\inserted{(L'objet ``le plus fin'' est du $C_\bullet(\mathscr X_\ast,\mathbb Z)$, on a alors $C_\bullet(\mathscr X_\ast,k)\simeq C_\bullet(\mathscr X_\ast,\mathbb Z)\otimes_{\mathbb Z} k$)}

Il est possible de \deleted{\unsure{au}} de définir une cobigèbre \inserted{diff} $C_{\bullet\bullet}^{\mathrm{DRpd}}(\mathscr X_\ast,\deleted{\unsure{\mathbb Z|k}}k)$ \inserted{$k$-coaugmentée} à copuiss. div, telle que l'on ait, pour tte $k^{+}$-alg. $k'$
\[
C^{\bullet\bullet}_{\mathrm{DRpd}}(\mathscr X_\ast,k')\simeq\mathrm{Hom}_k(C_{\bullet\bullet}^{\mathrm{DRpd}}(\mathscr X_\ast,k),k')
\]
(compatible avec ttes les structures).

On aura d'ailleurs
\[
C_{\bullet\bullet}^{\mathrm{DRpd}}(\mathscr X_\ast,k')\simeq C_{\bullet\bullet}^{\mathrm{DRpd}}(\mathscr X_\ast,k)\otimes_k k'
\]
L'objet le plus fin est $C_{\bullet\bullet}^{\mathrm{DRpd}}(\mathscr X_\ast,\mathbb Z)$. C'est lui qu'il conviendrait de considérer (au lieu de son ``dual'' $C^{\bullet\bullet}_{\mathrm{DRpd}}(\mathscr X_\ast,\mathbb Z)$) si on veut aborder b) c) d) sans conditions de finitude.

Notons que [tt comme $C^{\bullet\unsure{\ast|\bullet}}_{\mathrm{DRpd}}(\mathscr X_\ast,k)$, pour $\mathscr X_\ast$ variable, transforme $\underset{\longrightarrow}{\lim}$ quelq. en $\underset{\longleftarrow}{\lim}$], $C_{\bullet\bullet}^{\mathrm{DRpd}}(\mathscr X_\ast,\mathbb Z)$ \deleted{transforme} \inserted{commute aux} $\underset{\longrightarrow}{\lim}$ quelconques (pas néc. filtrantes)

% archive sheet number: 14
\scanpage{16}
\mspage{14}
% slanted
\margin{après f) ci-dessous}
\deleted{f}g) \uline{Faisceautisation} Il y a une définition évidente des complexes de De Rham à p.d. sur $k$ si $k$ est un Anneau comm. d'un topos. On aimerait, p. ex. en comparant un tel complexe à un autre qu. iso. dont les composantes soient flasques (mais y en a-t-il tjrs ?) définir des \unsure{op.s|op.} $\mathbb R f_\ast$ dans des cat. dérivées convenables pour de tels complexes, quand $f:\unsure{\mathscr X|X|\mathcal X}\to\mathscr Y$ est un morphisme de topos. \inserted{(approprié au besoin de \deleted{\unsure{dim}} coh. finie \ldots)} Si $k$ est une $\mathbb Q$-Algèbre, le $=$ Pb se rencontre d'ailleurs déjà pour les complexes du type De Rham-Sullivan, (et le Pb est ouvert -- et posé par Deligne -- quand $\mathscr X=(\unsure{\mathrm{Ens}|\mathscr Ens})^{\unsure{1|\Delta}}=$ topos des ens. cosimpliciaux, $\mathscr Y=$ topos ponctuel. Mais sauf erreur (si souvenirs sont exacts) il y a \unsure{un|une} \unsure{topos|topis|topo} qui marche, pour les espaces \inserted{topologiques} paracompacts \ldots
% archive sheet number: 15
